\documentclass[10pt,a4paper]{article}
\usepackage[utf8]{inputenc}
\usepackage[italian]{babel}
\usepackage{subfiles}
\usepackage{titling}
\usepackage{graphicx}
\usepackage{epigraph}
\usepackage[bottom]{footmisc}
\renewcommand*{\thefootnote}{[\arabic{footnote}]}
\usepackage{hyperref}
\hypersetup{
	colorlinks,
	citecolor=black,
	filecolor=black,
	linkcolor=black,
	urlcolor=black
}
\newcommand{\URI}[3][blue]{\href{#2}{\color{#1}{#3}}}%

\setlength{\droptitle}{-10em}
\pretitle{%
	\begin{center}
		\LARGE
		\centering
		\includegraphics[width=10cm,height=10cm]{unipd}\\[\bigskipamount]
	}
\posttitle{
	\includegraphics[width=350px,height=140px]{logo}\\[\bigskipamount]
	\includegraphics[width=2cm,height=2cm]{logo2}\\[\bigskipamount]
\end{center}
}
\title{
  Eatbrain.net - Analisi Usabilità \\
  Università degli Studi di Padova \\
  Web Information Management
}
\author{
  Gianluca Marraffa, 1121080
}
\date{Giugno 2018}


\begin{document}

\maketitle
\newpage
\pagenumbering{roman}
\tableofcontents
\newpage
\pagenumbering{arabic}

\section{Introduzione e scopo del documento}
Al giorno d'oggi qualsiasi associazione piccola o grande dispone di un sito web per presentare la propria immagine al mondo. Lo sviluppo di un portale personale, spesso integrato con piattaforme sociali (i social network), è diventato infatti imperativo per promuovere in modo efficace la propria ragione d'essere. Col tempo le tecnologie si sono raffinate e da semplici documenti di ipertesto, le pagine web sono diventate delle vere e proprie applicazioni interattive, solo un fattore però era, è e sarà per sempre una costante: l'usabilità.
\setlength{\epigraphwidth}{0.8\textwidth}
\epigraph{Usability and user experience design is about designing products to be effective, efficient, and satisfying. Specifically, ISO defines usability as the “extent to which a product can be used by specified users to achieve specified goals effectively, efficiently and with satisfaction in a specified context of use”}{ISO 9241-11}
Questo documento si impegna ad analizzare l'usabilità desktop del sito web
\begin{center}
	\URI{https://eatbrain.net/}{eatbrain.net}
\end{center}
 evidenziandone pregi e difetti per poi concludere con una valutazione generale della fruizione dei contenuti.
\newpage

\section{Analisi Preliminare}
\subfile{parts/preliminare}
\section{Analisi Complessiva}
\subfile{parts/complessiva}
\section{Home Page}
\section{Artists}
\section{Releases/Podcast}
\section{Events}
\section{Altre Pagine}
Contacts poca roba da valutare
Shop porta ad un altro dominio
\section{Valutazione e Considerazioni Finali}
Voto Dominio 9.5
Voto SEO 6
\appendix
\section{Tabella delle figure}

\end{document}